\section{Competition Retrospective}
\label{sec:competition-retrospective}

\subsection{Competition Setup}

\subsubsection{Timeline}

The \pokeagent Challenge ran from July 2025 through November 2025, culminating in presentations at NeurIPS 2025 in December. The timeline proceeded as follows: July launch with open ladder access, September hackathon for community building and technical talks, October qualifying rounds requiring 250+ games for statistical validity, and November finals featuring best-of-99 head-to-head matches. This structure allowed iterative improvement while maintaining evaluation integrity through held-out final assessments. The Speedrunning Track accepted submissions throughout with periodic leaderboard updates.

\subsubsection{Resources}

To support participants, we deployed several resources:
\begin{itemize}
    \item A dedicated custom Pokémon Showdown server hosting baseline agents and participant battles
    \item A RAG-powered Discord chatbot (``@pokeagent'') with a Professor Pokémon persona, using an embedding model and retrieval pipeline over organizer-curated documentation to answer participant questions
    \item Google Cloud Platform credits including Gemini API access
    \item A September hackathon with research talks streamed on YouTube
\end{itemize}

\subsection{Organizational Outcomes}

\paragraph{Participation Statistics}
The \pokeagent Challenge exceeded participation expectations:
\begin{itemize}
    \item \textbf{100+ active teams} with registered submissions across both tracks
    \item \textbf{650+ Discord community members} engaging in technical discussions
    \item \textbf{100K+ battles} on the competition Showdown server
    \item \textbf{22 valid Speedrunning Track submissions}, with 6 achieving 100\% completion
\end{itemize}

\paragraph{Dual Submission Tracks}
The dual-track structure successfully attracted participants from both the RL and LLM research communities. The Battling Track appealed to game AI and multi-agent learning researchers, while the Speedrunning Track attracted the long-context reasoning and language agent communities. Several teams competed in both tracks, developing unified architectures that could handle both strategic combat and exploration.

\paragraph{Cross-Platform Engagement}
The competition benefited significantly from the broader Pokémon AI zeitgeist of 2025. The Claude Plays Pokémon and Gemini Plays Pokémon communities provided a natural audience, and cross-promotion through Discord, Reddit (r/ClaudePlaysPokemon), and social media drove significant participation. The hackathon research talks, streamed on YouTube, attracted viewers from both academic and hobbyist communities.

\paragraph{Community Infrastructure}
Our RAG-powered ``@pokeagent'' Discord bot proved valuable for participant support, answering common questions about environment setup, baseline usage, and submission procedures using organizer-curated documentation. This reduced organizer burden while providing 24/7 assistance.

\paragraph{Timing Recommendations}
While NeurIPS encourages competitions to launch early, we found that peak participation and sustained engagement coincided with a 2--3 month window aligned with the Fall university semester (September--November). This timing particularly incentivized student involvement, as participants could integrate the competition into course projects and independent studies, creating a natural alignment between academic schedules and competition milestones.

\paragraph{Organizer Disclosure.} The organizing team includes the developers of PokéChamp and Metamon, which serve as baselines. To ensure fairness, organizer baselines were excluded from prize consideration, and several strong baselines were withheld during the competition to ensure that qualifying required genuine technical contribution beyond fine-tuning provided checkpoints. Of the 16 qualifying tournament spots, 13 were secured by teams that extended our public RL baselines, while 3 used independent approaches---reflecting that the released baselines provided an effective starting point rather than an unfair advantage. All evaluation was conducted on a shared server with identical conditions for all participants.
